% ========================================
% DANH MỤC KÝ HIỆU VÀ CHỮ VIẾT TẮT
% ========================================
\begin{center}
{\fontsize{18pt}{22pt}\selectfont\bfseries DANH MỤC KÝ HIỆU VÀ CHỮ VIẾT TẮT}
\end{center}
\phantomsection \addcontentsline{toc}{chapter}{DANH MỤC KÝ HIỆU VÀ CHỮ VIẾT TẮT}
\label{sec:acronyms}
\vspace{0.5cm}

\small
\begin{longtable}{|l|p{12cm}|}
\hline
\textbf{Viết tắt} & \textbf{Ý nghĩa} \\ \hline
\endfirsthead
\hline
\textbf{Viết tắt} & \textbf{Ý nghĩa} \\ \hline
\endhead
\hline
\endfoot
ACK & Acknowledgement (Bản tin xác nhận) \\ \hline
ACL & Access Control List (Danh sách kiểm soát truy cập) \\ \hline
AES & Advanced Encryption Standard (Tiêu chuẩn mã hóa nâng cao) \\ \hline
AF & Application Framework (Khung ứng dụng) \\ \hline
API & Application Programming Interface (Giao diện lập trình ứng dụng) \\ \hline
APS & Application Support Sub-layer (Phân lớp hỗ trợ ứng dụng trong Zigbee) \\ \hline
ASH & Asynchronous Serial Host (Giao thức kết nối Host - NCP không đồng bộ) \\ \hline
AWS & Amazon Web Services (Dịch vụ điện toán đám mây Amazon) \\ \hline
BLE & Bluetooth Low Energy (Công nghệ truyền thông Bluetooth năng lượng thấp) \\ \hline
CA & Certificate Authority (Tổ chức cấp phát chứng chỉ số) \\ \hline
CID & Correlation ID (Mã định danh liên kết) \\ \hline
CLI & Command Line Interface (Giao diện dòng lệnh) \\ \hline
CMD & Command (Lệnh điều khiển) \\ \hline
CMSIS & Cortex Microcontroller Software Interface Standard (Tiêu chuẩn giao tiếp phần mềm vi điều khiển Cortex) \\ \hline
CPU & Central Processing Unit (Bộ vi xử lý trung tâm) \\ \hline
CRC & Cyclic Redundancy Check (Mã kiểm tra thặng dư chu kỳ) \\ \hline
CRLF & Carriage Return Line Feed (Ký tự xuống dòng) \\ \hline
CRUD & Create, Read, Update, Delete (Các thao tác dữ liệu: Thêm, Đọc, Sửa, Xóa) \\ \hline
CSA & Connectivity Standards Alliance (Liên minh Tiêu chuẩn Kết nối) \\ \hline
CSDL & Cơ sở dữ liệu \\ \hline
CSMA/CA & Carrier Sense Multiple Access with Collision Avoidance (Đa truy cập nhận biết sóng mang tránh va chạm) \\ \hline
CTS & Clear To Send (Tín hiệu sẵn sàng nhận trong truyền thông nối tiếp) \\ \hline
DB & Database (Cơ sở dữ liệu) \\ \hline
DHT11 & Cảm biến nhiệt độ và độ ẩm kỹ thuật số tích hợp \\ \hline
DMA & Direct Memory Access (Truy cập bộ nhớ trực tiếp) \\ \hline
E2E & End-to-End (Toàn trình / Từ đầu cuối đến đầu cuối) \\ \hline
EC2 & Elastic Compute Cloud (Dịch vụ máy chủ ảo đám mây của AWS) \\ \hline
EUI / EUI64 & Extended Unique Identifier (Định danh mở rộng duy nhất 64-bit của thiết bị Zigbee) \\ \hline
EZSP & EmberZNet Serial Protocol (Giao thức truyền thông nối tiếp EmberZNet) \\ \hline
FFD & Full Function Device (Thiết bị đầy đủ chức năng trong mạng Zigbee) \\ \hline
GBL & Gecko Bootloader (Tệp tin phân phối firmware dạng bootloader của Silicon Labs) \\ \hline
GPIO & General Purpose Input/Output (Cổng vào ra chung) \\ \hline
GSDK & Gecko SDK (Bộ phát triển phần mềm Gecko) \\ \hline
GUI & Graphical User Interface (Giao diện người dùng đồ họa) \\ \hline
HTTP & Hypertext Transfer Protocol (Giao thức truyền tải siêu văn bản) \\ \hline
HTTPS & Hypertext Transfer Protocol Secure (Giao thức truyền tải siêu văn bản bảo mật) \\ \hline
I2C & Inter-Integrated Circuit (Giao thức truyền thông nối tiếp giữa các vi mạch) \\ \hline
IC & Install Code (Mã cài đặt bảo mật thiết bị trong Zigbee) \\ \hline
ID & Identifier (Mã định danh) \\ \hline
IDE & Integrated Development Environment (Môi trường phát triển tích hợp) \\ \hline
IEEE & Institute of Electrical and Electronics Engineers (Viện kỹ sư Điện và Điện tử) \\ \hline
IP & Internet Protocol (Giao thức Internet) \\ \hline
IPC & Inter-Process Communication (Truyền thông giữa các tiến trình) \\ \hline
IoT & Internet of Things (Mạng lưới vạn vật kết nối Internet) \\ \hline
JSON & JavaScript Object Notation (Định dạng trao đổi dữ liệu dạng văn bản nhẹ) \\ \hline
JWT & JSON Web Token (Mã xác thực dạng chuỗi ký tự JSON mã hóa) \\ \hline
LCD & Liquid Crystal Display (Màn hình tinh thể lỏng) \\ \hline
LED & Light Emitting Diode (Điốt phát quang) \\ \hline
LQI & Link Quality Indicator (Chỉ số chất lượng liên kết vật lý) \\ \hline
LWT & Last Will and Testament (Cơ chế di chúc trong MQTT) \\ \hline
MAC & Media Access Control (Phân lớp điều khiển truy cập môi trường truyền dẫn) \\ \hline
MCU & Microcontroller Unit (Khối vi điều khiển) \\ \hline
MQTT & Message Queuing Telemetry Transport (Giao thức truyền tải tin nhắn dạng publish/subscribe) \\ \hline
NAK & Negative Acknowledgement (Bản tin báo nhận lỗi truyền tin) \\ \hline
NCP & Network Co-Processor (Bộ vi xử lý mạng phụ trợ) \\ \hline
NVM3 & Non-Volatile Memory 3rd Generation (Hệ thống lưu trữ bộ nhớ không bay hơi thế hệ thứ 3 của Silicon Labs) \\ \hline
NWK & Network Layer (Lớp mạng trong mô hình Zigbee) \\ \hline
ORM & Object-Relational Mapping (Ánh xạ quan hệ - đối tượng) \\ \hline
OS & Operating System (Hệ điều hành) \\ \hline
OTA & Over-the-Air (Công nghệ nâng cấp firmware không dây) \\ \hline
PAN & Personal Area Network (Mạng khu vực cá nhân) \\ \hline
PHY & Physical Layer (Lớp vật lý) \\ \hline
PIR & Passive Infrared Sensor (Cảm biến hồng ngoại thụ động) \\ \hline
PSA & Platform Security Architecture (Kiến trúc bảo mật nền tảng) \\ \hline
PTI & Packet Trace Interface (Giao diện giám sát gói tin phần cứng) \\ \hline
QoS & Quality of Service (Chất lượng dịch vụ trong truyền tin) \\ \hline
RAIL & Radio Abstraction Interface Layer (Lớp giao diện trừu tượng hóa sóng vô tuyến của Silicon Labs) \\ \hline
RAM & Random Access Memory (Bộ nhớ truy cập ngẫu nhiên) \\ \hline
RBAC & Role-Based Access Control (Kiểm soát truy cập dựa trên vai trò) \\ \hline
REST & Representational State Transfer (Kiểu kiến trúc truyền trạng thái đại diện để thiết kế API Web) \\ \hline
RF & Radio Frequency (Tần số sóng vô tuyến) \\ \hline
RFD & Reduced Function Device (Thiết bị giảm thiểu chức năng trong mạng Zigbee) \\ \hline
RPi & Raspberry Pi (Máy tính nhúng bo mạch đơn phổ biến) \\ \hline
RTOS & Real-Time Operating System (Hệ điều hành thời gian thực) \\ \hline
RTS & Request To Send (Tín hiệu yêu cầu gửi dữ liệu trong UART) \\ \hline
RX & Receive (Tín hiệu nhận dữ liệu nối tiếp) \\ \hline
SDK & Software Development Kit (Bộ công cụ phát triển phần mềm) \\ \hline
SED & Sleepy End Device (Thiết bị đầu cuối Zigbee chế độ ngủ tiết kiệm pin) \\ \hline
SHA256 & Secure Hash Algorithm 256-bit (Thuật toán băm bảo mật 256-bit) \\ \hline
SoC & System on Chip (Hệ thống trên một vi mạch) \\ \hline
SPI & Serial Peripheral Interface (Giao thức giao tiếp ngoại vi nối tiếp đồng bộ) \\ \hline
SSE & Server-Sent Events (Công nghệ đẩy dữ liệu một chiều từ máy chủ xuống trình duyệt) \\ \hline
SSv5 & Simplicity Studio version 5 (Môi trường phát triển phần mềm của Silicon Labs phiên bản 5) \\ \hline
TCP & Transmission Control Protocol (Giao thức điều khiển truyền vận) \\ \hline
TCP/IP & Transmission Control Protocol / Internet Protocol (Bộ giao thức truyền thông Internet) \\ \hline
TLS & Transport Layer Security (Bảo mật tầng truyền tải) \\ \hline
TX & Transmit (Tín hiệu truyền dữ liệu nối tiếp) \\ \hline
UART & Universal Asynchronous Receiver-Transmitter (Bộ truyền nhận nối tiếp không đồng bộ) \\ \hline
UI & User Interface (Giao diện người dùng) \\ \hline
USB & Universal Serial Bus (Chuẩn giao tiếp bus nối tiếp vạn năng) \\ \hline
UUID & Universally Unique Identifier (Mã định danh duy nhất toàn cầu) \\ \hline
UX & User Experience (Trải nghiệm người dùng) \\ \hline
VCOM & Virtual COM Port (Cổng COM ảo) \\ \hline
WSTK & Wireless Starter Kit (Bộ kit phát triển phần cứng không dây của Silicon Labs) \\ \hline
XOFF & Transmitter Off (Ký tự báo ngừng truyền dữ liệu trong điều khiển luồng phần mềm) \\ \hline
XON & Transmitter On (Ký tự báo tiếp tục truyền dữ liệu trong điều khiển luồng phần mềm) \\ \hline
XOR & Exclusive OR (Phép toán logic OR loại trừ) \\ \hline
ZC & Zigbee Coordinator (Thiết bị điều phối mạng Zigbee) \\ \hline
ZCL & Zigbee Cluster Library (Thư viện tập hợp các thuộc tính và lệnh chuẩn của Zigbee) \\ \hline
ZDO & Zigbee Device Object (Đối tượng thiết bị quản lý cấu hình lớp mạng Zigbee) \\ \hline
ZED & Zigbee End Device (Thiết bị đầu cuối mạng Zigbee) \\ \hline
ZR & Zigbee Router (Thiết bị định tuyến mạng Zigbee) \\ \hline
\end{longtable}
\normalsize

\newpage
